\section{The initial ensemble and the Gaussian target}\label{sec:initial-target}

This section proves two facts independently of the dynamical expansion:
the initial fluctuation theorem, and well-posedness of the Gaussian law
appearing in \Cref{thm:main}.  It also records the initial correlation
bounds used in the expansion.

\subsection{Initial factorial measures and connected graphs}

Let $F_{h,0}^{\eps}$ be the density of the ordered factorial measure
divided by $\mu^h$.  Thus, for a nonnegative measurable $\Phi$,
\begin{equation}\label{eq:initial-factorial-definition}
 \int F_{h,0}^{\eps}(Z_h)\Phi(Z_h)\,\dd Z_h
 =\mu^{-h}\E\sum_{i_1,\ldots,i_h\text{ distinct}}
       \Phi(z_{i_1},\ldots,z_{i_h}).
\end{equation}
Write $[h]=\{1,\ldots,h\}$, let $\mathcal P(V)$ denote the set of
partitions of a finite set $V$, and set $Z_B=(z_j)_{j\in B}$.
The normalized factorial cumulant densities, also called connected
densities, $c_{h,0}^{\eps}$ are defined by
\begin{equation}\label{eq:initial-partition-definition}
 F_{h,0}^{\eps}(Z_h)
   =\sum_{\pi\in\mathcal P([h])}
       \prod_{B\in\pi}c_{|B|,0}^{\eps}(Z_B).
\end{equation}
In particular, $g_0=c_{1,0}^{\eps}=F_{1,0}^{\eps}$.
Set
\[
 \varrho_0(x)=\int_{\R^3}f_0(x,v)\,\dd v,
 \qquad b_\eps=\frac{4\pi}{3}\|\varrho_0\|_\infty\eps^3,
 \qquad d_\eps=\mu b_\eps.
\]
The initial assumptions imply that $\varrho_0$ and its spatial gradient
are bounded, and $d_\eps=O(\eps)$.
We write $(N)_p=N(N-1)\cdots(N-p+1)$ for the falling factorial,
with $(N)_0=1$.

\begin{proposition}[Initial ensemble]\label{prop:initial}
For all sufficiently small $\eps$, the following assertions hold.
\begin{enumerate}[label=\textup{(\roman*)}]
\item The factorial densities satisfy
\begin{equation}\label{eq:initial-factorial-upper}
 0\le F_{h,0}^{\eps}\le\prod_{j=1}^h f_0(z_j),
 \qquad (1-C\eps)f_0\le g_0\le f_0.
\end{equation}
Consequently $\|g_0-f_0\|_{\Bol,2\beta}=O(\eps)$ and, for every
fixed integer $p\ge1$,
\begin{equation}\label{eq:initial-N-moments}
 \E(N)_p\le\mu^p,
 \qquad \sup_{0<\eps\le\eps_0}\E(N/\mu)^p<\infty.
\end{equation}
\item With $\upsilon=1/81$, for every $h\ge2$,
\begin{equation}\label{eq:initial-coarse-links}
 |c_{h,0}^{\eps}(Z_h)|
 \le (Ch)^{Ch}\prod_{j=1}^h f_0(z_j)
   \sum_{\mathcal T\text{ tree on }[h]}
     \prod_{\{i,j\}\in\mathcal T}
       \bigl(\ind_{|x_i-x_j|\le\eps}+\eps^\upsilon\bigr).
\end{equation}
Here $C$ and the upper bound on $\eps$ are independent of $h$.
For fixed $h\ge2$ and $D>0$ there is $C_{D,h}<\infty$ such that
\begin{equation}\label{eq:initial-sharp-links}
 |c_{h,0}^{\eps}(Z_h)|
 \le C_{D,h}\prod_{j=1}^h f_0(z_j)
   \sum_{\mathcal T\text{ tree on }[h]}
     \prod_{\{i,j\}\in\mathcal T}
       \bigl(\ind_{|x_i-x_j|\le C_{D,h}\eps}+\eps^D\bigr).
\end{equation}
Moreover,
\begin{equation}\label{eq:initial-integrated-cumulants}
 \|c_{h,0}^{\eps}\|_{L^1}\le C_h\eps^{3(h-1)}
 \qquad(h\ge2).
\end{equation}
\item For every finite list of bounded real tests, the exact-centered
initial fluctuations converge jointly to a centered Gaussian vector
with covariance $\int f_0\phi\psi$.
For a smooth compactly supported real test $\phi$, one also has
\begin{equation}\label{eq:initial-centering-shift}
 \lim_{\eps\to0}\sqrt\mu
 \left(\E\pi_0^\eps(\phi)-\int f_0\phi\right)
 =-\frac{4\pi}{3}\int_{\R^6}
       f_0(x,v)\phi(x,v)\varrho_0(x)\,\dd x\,\dd v.
\end{equation}
\end{enumerate}
\end{proposition}

\begin{proof}
For fixed $Z_h$, let $\mathcal Z_\eps[Z_h]$ be the partition function
of the remaining particles with their usual mutual exclusion and with
the additional restriction that they avoid all the balls
$B(x_j,\eps)$.  Directly from the ensemble,
\begin{equation}\label{eq:initial-insertion}
 F_{h,0}^{\eps}(Z_h)=
 \ind_{\mathcal D_h^\eps}(Z_h)\prod_{j=1}^h f_0(z_j)
 \frac{\mathcal Z_\eps[Z_h]}{\mathcal Z_\eps}.
\end{equation}
The ratio is at most one, proving the factorial upper bound.  When
$h=1$, it is the probability under the original ensemble that the ball
$B(x_1,\eps)$ contains no particle center.  The factorial upper bound
with one particle and the union bound show that this probability is at
least $1-\mu b_\eps$.  This proves \eqref{eq:initial-factorial-upper}.
Integrating the factorial upper bound gives the first part of
\eqref{eq:initial-N-moments}; expressing $N^p$ as a finite positive
linear combination of $(N)_j$, $j\le p$, gives its second part.

We give the connected expansion and its convergence explicitly.
Write $\chi_{ij}=\ind_{|x_i-x_j|\le\eps}$ and set
\[
 \mathcal U_1=1,
 \qquad
 \mathcal U_n(Z_n)=
   \sum_{H\text{ connected on }[n]}
      \prod_{\{i,j\}\in H}(-\chi_{ij})\quad(n\ge2).
\]
The Penrose partition scheme~\cite{Penrose1967}, in the form proved
in~\cite[Proposition~2.3.3]{BGSS2023}, gives the following tree bound.
\begin{equation}\label{eq:initial-tree-inequality}
 |\mathcal U_n(Z_n)|\le
   \sum_{\mathcal T\text{ tree on }[n]}
       \prod_{\{i,j\}\in\mathcal T}\chi_{ij}.
\end{equation}
To see this, root a connected labelled graph at vertex $1$, assign its
graph-distance levels, and choose the least-labelled neighbor in the
preceding level as the parent of each other vertex.  For a fixed
resulting tree $\mathcal T$, the graphs that give this tree form the
interval $\mathcal T\subset H\subset R(\mathcal T)$.  Here
$R(\mathcal T)$ contains all same-level edges and every edge from a
vertex to a preceding-level vertex whose label is at least the label
of its chosen parent.  No other edge is permitted.  Summing this
interval yields
\[
 \prod_{e\in\mathcal T}(-\chi_e)
 \prod_{e\in R(\mathcal T)\setminus\mathcal T}(1-\chi_e).
\]
These intervals partition all connected graphs, and every factor in
the second product lies in $[0,1]$, proving
\eqref{eq:initial-tree-inequality}.

Integrating a tree successively from its leaves, and keeping one free
vertex, gives
\begin{equation}\label{eq:initial-unrooted-tree-integral}
 \int f_0^{\otimes n}|\mathcal U_n|
 \le n^{n-2}b_\eps^{n-1}\qquad(n\ge2).
\end{equation}
For a bounded measurable activity multiplier $a$, the formal
connected-graph identity is therefore an absolutely convergent
identity whenever $e\mu b_\eps\|a\|_\infty<1$:
\begin{equation}\label{eq:initial-log-partition}
 \log\mathcal Z_\eps[a]
  =\sum_{n\ge1}\frac{\mu^n}{n!}
    \int\mathcal U_n(Z_n)\prod_{j=1}^n f_0(z_j)a(z_j)\,\dd Z_n.
\end{equation}
For completeness, first multiply the activity by a complex scalar
$w$.  The identity follows coefficient by coefficient near $w=0$
from the finite exponential formula for connected components.  The
tree bound makes the right side absolutely convergent on the stated
disk, so exponentiation and analytic continuation prove the identity
there.  This also fixes the branch of the logarithm.  The same bounds
justify every fixed number of derivatives in bounded source
multipliers.

Apply \eqref{eq:initial-log-partition} to $a=1+u$.  The logarithm of
$\E\prod_i(1+u(z_i))$ has coefficients $\mu^h c_{h,0}^{\eps}/h!$ by
\eqref{eq:initial-partition-definition}.  Differentiating the
absolutely convergent series gives
\begin{equation}\label{eq:initial-rooted-mayer}
 c_{h,0}^{\eps}(Z_h)
 =\prod_{j=1}^h f_0(z_j)
   \sum_{n\ge0}\frac{\mu^n}{n!}
      \int f_0^{\otimes n}(Z_{h+1:h+n})
         \mathcal U_{h+n}(Z_{h+n})\,\dd Z_{h+1:h+n}.
\end{equation}
We next give bounds uniform in the number $h$ of fixed roots.
From any tree on the $h+n$ vertices, delete $h-1$ edges so that each
remaining component contains exactly one fixed root.  Such a deletion
is obtained by repeatedly cutting an edge on a path between two roots.
Dropping the deleted constraints and integrating the auxiliary leaves
then costs at most $b_\eps^n$.  For $n\ge1$,
\begin{equation}\label{eq:initial-rooted-count}
 \frac{(h+n)^{h+n-2}}{n!}
 \le (eh)^h e^{2n}.
\end{equation}
Indeed, $n!\ge(n/e)^n$,
$(1+h/n)^n\le e^h$, and $(h+n)^h\le h^h e^n$.
The $n=0$ count is also bounded by $(Ch)^{Ch}$.
It follows that \eqref{eq:initial-rooted-mayer} is absolutely convergent
for $\eps\le\eps_0$, with this $\eps_0$ independent of $h$.

Let $q$ be the number of connected components of the graph on the
fixed roots with edges $|x_i-x_j|\le\eps$.  If $q\ge2$, every
contributing full tree must attach each such component to an auxiliary
vertex.  An auxiliary vertex can be adjacent to at most $27$ different
root components: choose one adjacent root from each; their mutual
distances exceed $\eps$, and their disjoint balls of radius $\eps/2$
lie in a ball of radius $3\eps/2$.  Thus a contributing tree with
$n$ auxiliary vertices has $q\le27n$.
Using \eqref{eq:initial-rooted-count} and $d_\eps=O(\eps)$, we obtain
\[
 \frac{|c_{h,0}^{\eps}(Z_h)|}{\prod_j f_0(z_j)}
 \le (Ch)^{Ch}\eps^{q/27}
 \le (Ch)^{Ch}\eps^{(q-1)/81}\qquad(q\ge2),
\]
with the quotient interpreted through
\eqref{eq:initial-rooted-mayer} when some $f_0(z_j)=0$.
For $q=1$ the same argument without the restriction on $n$ gives
$(Ch)^{Ch}$.  There exists a tree of the roots using edges inside each
of the $q$ components and exactly $q-1$ edges between components.
Its product in \eqref{eq:initial-coarse-links} is at least
$\eps^{(q-1)/81}$.  This proves the coarse link estimate.

For fixed $h,D$, truncate \eqref{eq:initial-rooted-mayer} at an integer
$n_0$ with $n_0+1\ge D(h-1)$.  Its tail is bounded by
\[
 C_{D,h}\eps^{D(h-1)}\prod_j f_0(z_j).
\]
In any contributing term
with $n\le n_0$, all fixed roots are connected by paths of at most
$h+n_0-1$ links of length at most $\eps$.  Thus their diameter is at
most $(h+n_0)\eps$.  On that set the tree sum in
\eqref{eq:initial-sharp-links} is at least one when $C_{D,h}$ is chosen
large enough; everywhere it is at least $\eps^{D(h-1)}$.
This proves \eqref{eq:initial-sharp-links}.  Integrating all roots in
\eqref{eq:initial-rooted-mayer} and using
\eqref{eq:initial-unrooted-tree-integral} instead gives
\[
 \|c_{h,0}^{\eps}\|_1
 \le b_\eps^{h-1}
       \sum_{n\ge0}\frac{(h+n)^{h+n-2}}{n!}d_\eps^n
 \le C_h\eps^{3(h-1)},
\]
as asserted.

For a real random variable $X$ with exponential moments near the
origin, write
\[
 \kappa_k(X)=\left.\frac{\dd^k}{\dd u^k}\log\E e^{uX}\right|_{u=0},
 \qquad k\ge1,
\]
for its ordinary cumulants.
To prove the initial central limit theorem, let $\phi$ be bounded and
real and put $S_\phi=\sum_i\phi(z_i)$.  Substitute
$a(z)=e^{u\phi(z)}$ into \eqref{eq:initial-log-partition} and
differentiate $k$ times at $u=0$.  The resulting ordinary cumulant is
\[
 \kappa_k(S_\phi)=
 \sum_{n\ge1}\frac{\mu^n}{n!}
  \int f_0^{\otimes n}\mathcal U_n
        \left(\sum_{j=1}^n\phi(z_j)\right)^k.
\]
For each fixed $k$, \eqref{eq:initial-unrooted-tree-integral} gives
\begin{equation}\label{eq:initial-ordinary-cumulants}
 \kappa_k(S_\phi)=\mu\int f_0\phi^k+O_{k,\phi}(\mu\eps).
\end{equation}
In fact the absolute error is bounded by
$\mu\|\phi\|_\infty^k\sum_{n\ge2}
n^{n+k-2}d_\eps^{n-1}/n!$, which is $O_{k,\phi}(\mu\eps)$.
After exact centering and division by $\sqrt\mu$, the second cumulant
tends to $\int f_0\phi^2$ and every higher cumulant tends to zero.
The moment--cumulant identity gives every Gaussian moment.  To see
directly that this implies scalar convergence, compare the
characteristic functions of the normalized variable and the Gaussian
with variance \(\sigma^2=\int f_0\phi^2\), using their Taylor
polynomials of degree \(2n-1\).  After taking \(\limsup_{\eps\to0}\),
the sum of their two remainder bounds is at most
\(2|u|^{2n}\sigma^{2n}/(2^n n!)\), which tends to zero for each
real \(u\).  Applying scalar convergence to each linear combination
of the tests gives joint convergence.  This argument also includes
zero variance.

Finally, the $n=1$ auxiliary term of
\eqref{eq:initial-rooted-mayer} with $h=1$ gives, uniformly in $z$,
\[
 g_0(z)=f_0(z)\left[
  1-\mu\int_{|x-x_*|\le\eps}\varrho_0(x_*)\,\dd x_*
  +O(\eps^2)\right].
\]
The bounded gradient of $\varrho_0$ implies
\[
 \int_{|x-x_*|\le\eps}\varrho_0(x_*)\,\dd x_*
 =\frac{4\pi}{3}\eps^3\varrho_0(x)+O(\eps^4).
\]
Since $\sqrt\mu=\eps^{-1}$, integration against $\phi$ proves
\eqref{eq:initial-centering-shift}.
\end{proof}

\begin{remark}
The shift in \eqref{eq:initial-centering-shift} is generally nonzero.
The exact microscopic centering in \Cref{thm:main} is therefore
essential even at the initial time.  On the other hand,
\eqref{eq:initial-integrated-cumulants} gives
$\mu\|c_{2,0}^{\eps}\|_1=O(\eps)$, so the limiting initial covariance
has only the same-particle contribution, $\int f_0\phi\psi$.
\end{remark}

\subsection{Velocity moments and the backward propagator}

Write $M=\sup_{t\le T,x,v}e^{2\beta|v|^2}f_t(x,v)$.
We first record the global integrability needed for the noise form;
it does not follow from a spatially uniform bound alone.

\begin{lemma}[Global weighted integrability]\label{lem:target-moments}
For every $0<q<2\beta$,
\begin{equation}\label{eq:target-velocity-moments}
 \sup_{0\le t\le T}\int_{\R^6}
       f_t(x,v)e^{q|v|^2}\,\dd x\,\dd v<\infty.
\end{equation}
\end{lemma}

\begin{proof}
Let $Q^+$ denote the positive gain part of $Q$, and define
$P_tg=Q^+(g,f_t)$, keeping the second factor fixed.
Write $\mathsf T_t g(x,v)=g(x-tv,v)$ for density transport, and
\[
 (\mathcal Vg)_t=\int_0^t\mathsf T_{t-s}P_sg_s\,\dd s.
\]
Positivity and the transport Duhamel formula imply
\begin{equation}\label{eq:gain-series-majorant}
 f_t\le\sum_{j=0}^{n-1}
       (\mathcal V^j\mathsf T_{\cdot}f_0)_t
       +(\mathcal V^nf)_t.
\end{equation}
We remove the remainder without presupposing spatial integrability.
On the Gaussian decay spaces
$Y_a=\{g:\|g\|_{Y_a}=\sup e^{a|v|^2}|g(x,v)|<\infty\}$,
energy conservation at a collision gives, for $0<b<a\le2\beta$,
\[
 |P_tg(x,v)|\le CM\|g\|_{Y_a}(1+|v|)e^{-a|v|^2},
 \qquad
 \|P_tg\|_{Y_b}
 \le CM\bigl(1+(a-b)^{-1/2}\bigr)\|g\|_{Y_a}.
\]
Constants are uniform when $b$ is bounded away from zero.  Distribute
the exponent interval from $\beta$ to $\beta/2$ equally among the
$n$ factors in the remainder.  Transport preserves these norms and
the time simplex has volume $T^n/n!$, so
\[
 \sup_{t\le T}\| (\mathcal V^nf)_t\|_{Y_{\beta/2}}
 \le M\frac{(CT)^n}{n!}
        \bigl(1+\sqrt{2n/\beta}\bigr)^n\longrightarrow0.
\]
Thus $f$ is bounded pointwise by the full positive gain series.

Let $X_q=L^1(e^{q|v|^2}\dd x\dd v)$.  The pre/postcollision
substitution, followed by energy conservation and the bound on $f_t$,
gives
\[
 \|P_tg\|_{X_q}
 \le CM\int |g(x,v)|(1+|v|)e^{q|v|^2}\,\dd x\,\dd v
 \le CM\bigl(1+\delta^{-1/2}\bigr)\|g\|_{X_{q+\delta}},
\]
uniformly for exponents in a compact subinterval of $(0,2\beta)$.
For example, in the first inequality use
$e^{q|v'|^2}\le e^{q|v|^2+q|v_*|^2}$ and integrate the second
factor against its Gaussian bound.  The initial Bol norm implies
$f_0\in X_{q_0}$ for every $q_0<2\beta$.  Given $q<q_0$, distribute
$q_0-q$ among the factors of the positive gain series.  The same
convergent bound with $n^{n/2}/n!$ proves that its $X_q$ norm is
uniformly finite for $t\le T$.  Integrating the pointwise majorant
proves \eqref{eq:target-velocity-moments}.
\end{proof}

For $a\ge0$ define the Banach test spaces
\begin{equation}\label{eq:target-test-spaces}
 E_a^0=\{h\in C(\R^6):e^{-a|v|^2}h(x,v)\in C_0(\R^6)\},
 \qquad \|h\|_a=\sup e^{-a|v|^2}|h(x,v)|.
\end{equation}
Smooth compactly supported tests are dense in each $E_a^0$.
Let $K_t=A_t-v\cdot\nabla_x$ and let
$\mathsf S_u h(x,v)=h(x+uv,v)$ denote free transport of tests.

\begin{lemma}[Backward evolution]\label{lem:target-propagator}
For $0\le a<b<2\beta$ and $0\le r\le t\le T$, there is a bounded
backward propagator $U(r,t):E_a^0\to E_b^0$.  It solves
\begin{equation}\label{eq:target-backward-mild}
 U(r,t)h=\mathsf S_{t-r}h
   +\int_r^t\mathsf S_{s-r}K_sU(s,t)h\,\dd s,
\end{equation}
has the composition law with intermediate exponents, and is the
unique mild propagator in these scales.  For $\phi\in C_c^\infty$,
\begin{equation}\label{eq:target-upper-end-identity}
 U(r,t)\phi-\phi=\int_r^t U(r,s)A_s\phi\,\dd s.
\end{equation}
\end{lemma}

\begin{proof}
For $a$ in a compact subinterval of $[0,2\beta)$, energy conservation
and the Gaussian bound give
\begin{equation}\label{eq:target-K-estimate}
 |K_th(x,v)|\le CM\|h\|_a(1+|v|)e^{a|v|^2},
 \qquad
 \|K_th\|_b\le CM\bigl(1+(b-a)^{-1/2}\bigr)\|h\|_a.
\end{equation}
Indeed $|\Delta_h|\le4\|h\|_a e^{a(|v|^2+|v_*|^2)}$ and
$\int B(v-v_*,\omega)\dd\omega=\pi|v-v_*|$.
The operator also maps $E_a^0$ into $E_b^0$.  To check the spatial
tail, replace $\|h\|_a$ in the first bound by
$\eta_h(x)=\sup_v e^{-a|v|^2}|h(x,v)|$, which tends to zero as
$|x|\to\infty$.  The exponent loss controls the velocity tail.
Classical continuity of $f$, local dominated convergence, and these
tail bounds show that $t\mapsto K_th$ is continuous in $E_b^0$.
Free transport is a strongly continuous isometry on each $E_a^0$.

Iterate \eqref{eq:target-backward-mild}.  Giving each of the $n$
collision factors the exponent allowance $(b-a)/n$ bounds the
$n$th time-ordered term by
\begin{equation}\label{eq:target-Dyson-bound}
 \frac{(CM(t-r))^n}{n!}
       \bigl(1+\sqrt{n/(b-a)}\bigr)^n\|h\|_a.
\end{equation}
The series converges for every finite $T$.  Its terms are Bochner
integrals on the stated spaces.  Splitting a time-ordered simplex at
an intermediate time proves composition; the same estimate applied
to an iterated homogeneous remainder proves uniqueness, leaving a
positive exponent allowance between its two endpoint spaces.

For a compactly supported smooth $\phi$, the function $A_s\phi$
has compact spatial support, grows at most linearly in $|v|$, and is
continuous in $s$ in every $E_a^0$ with $a>0$.  Differentiate the
series at its upper endpoint $t$.  Differentiating its final transport
factor gives $U(r,t)(v\cdot\nabla_x\phi)$, and the upper faces of
the simplices give $U(r,t)K_t\phi$.  Both differentiated series
converge with the bound \eqref{eq:target-Dyson-bound}, using a small
extra exponent allowance for $K_t\phi$.  Hence
$\partial_tU(r,t)\phi=U(r,t)A_t\phi$, which proves
\eqref{eq:target-upper-end-identity}.
\end{proof}

\subsection{Construction and uniqueness of the Gaussian field}

The backward covariance formulation below also appears in the short-time
fluctuation theory of~\cite[Section~6.1]{BGSS2023}.  Here the preceding
moment and propagator estimates construct the Gaussian field on the full
prescribed regular interval \([0,T]\).

\begin{proposition}[Gaussian target]\label{prop:gaussian-target}
The fluctuation equation in \Cref{thm:main} has a centered Gaussian
solution on the test scales $E_a^0$, $0\le a<\beta$.  It is a
continuous linear map from each such test space to $L^2$ of its
probability space, uniformly on $[0,T]$.  Its covariance is
\begin{equation}\label{eq:target-covariance}
 \begin{split}
 \E[\zeta_t(\phi)\zeta_s(\psi)]
 ={}&\int f_0(z)\,U(0,t)\phi(z)\,U(0,s)\psi(z)\,\dd z\\
 &+\int_0^{\min(s,t)}
     C_r\bigl(U(r,t)\phi,U(r,s)\psi\bigr)\,\dd r.
 \end{split}
\end{equation}
It is unique in the class of solutions extending to these scales
with uniformly bounded second-moment operator norms.  In particular,
the law required in \Cref{thm:main} is well defined.
\end{proposition}

\begin{proof}
For $b<\beta$, energy conservation and \Cref{lem:target-moments} give
\begin{equation}\label{eq:target-source-bound}
 |C_t(h,k)|\le CM\|h\|_b\|k\|_b
    \int f_t(x,v)(1+|v|)e^{2b|v|^2}\,\dd x\,\dd v
 \le C_b\|h\|_b\|k\|_b.
\end{equation}
Choose the moment exponent slightly above $2b$ to absorb the factor
$1+|v|$.  The initial quadratic form has the corresponding bound.

Let $I_0$ and $I$ be independent isonormal Gaussian maps on,
respectively, $L^2(f_0\dd z)$ and $L^2(\nu)$, where
\[
 \dd\nu=\frac12 B(v-v_*,\omega) f_r(x,v)f_r(x,v_*)
   \,\dd r\,\dd x\,\dd v\,\dd v_*\,\dd\omega.
\]
Such maps can be constructed by choosing orthonormal bases and
independent standard normal coefficients; their defining series
converge in $L^2$.  Define
\begin{equation}\label{eq:target-isonormal-formula}
 \zeta_t(h)=I_0\bigl(U(0,t)h\bigr)
       +I\bigl(\ind_{r\le t}\Delta_{U(r,t)h}\bigr),
 \qquad
 W_t(h)=I\bigl(\ind_{r\le t}\Delta_h\bigr).
\end{equation}
For $h\in E_a^0$, propagate to any exponent $b\in(a,\beta)$ and
apply \eqref{eq:target-source-bound}.  This proves that the maps are
well defined and that
\begin{equation}\label{eq:target-field-bound}
 \sup_{t\le T}\|\zeta_t(h)\|_{L^2(\Omega)}\le C_a\|h\|_a.
\end{equation}
They are jointly Gaussian and linear, and their covariance is
\eqref{eq:target-covariance}.  Orthogonality of disjoint time intervals
in $L^2(\nu)$ gives the independent increments of $W$; independence of
$I_0$ and $I$ gives the stated initial independence.

We verify the weak equation, including the possibly unbounded drift
test.  For $\phi\in C_c^\infty$, use
\eqref{eq:target-upper-end-identity} and the boundedness of the two
isonormal maps to interchange them with Bochner integrals.  All the
triangular kernels involved have finite squared norm by
\eqref{eq:target-source-bound}.  One obtains
\[
 \begin{split}
 \int_0^t\zeta_s(A_s\phi)\,\dd s
 ={}&I_0\bigl(U(0,t)\phi-\phi\bigr)\\
 &+I\bigl(\ind_{r\le t}
                \Delta_{U(r,t)\phi-\phi}\bigr).
 \end{split}
\]
Together with \eqref{eq:target-isonormal-formula}, this is precisely
\[
 \zeta_t(\phi)=\zeta_0(\phi)
       +\int_0^t\zeta_s(A_s\phi)\,\dd s+W_t(\phi).
\]

For uniqueness, work on the probability space of any candidate solution.
Construct the backward mild field there from its own initial field and
noise: the deterministic integrals are mean-square limits of compact
tests and time-step functions, with the bounds just proved.
Subtract this field from the candidate and call the difference $D_t$.
It has zero initial field and noise.  Its compact-test weak equation and
uniform second-moment bound give mean-square continuity.  Riemann sums
then permit the time-dependent tests $\mathsf S_{t-s}\phi$, whose
supports lie in one compact set for $0\le s\le t$.  They give
the transport mild identity
\[
 D_t(\phi)=\int_0^t
       D_s\bigl(K_s\mathsf S_{t-s}\phi\bigr)\,\dd s.
\]
It extends by density and \eqref{eq:target-K-estimate} to exponent
scales with a positive allowance.
For $0\le a<b<\beta$, iteration $n$ times and
\eqref{eq:target-field-bound} bound its $L^2$ norm by
\[
 C_b\|\phi\|_a\frac{(CT)^n}{n!}
             \bigl(1+\sqrt{n/(b-a)}\bigr)^n\longrightarrow0.
\]
This proves uniqueness in the stated class and therefore uniqueness
in law of the constructed Gaussian solution.

This construction is a Gaussian generalized field in the explicit
mean-square test-space sense of \eqref{eq:target-field-bound}; it does
not assert membership in the Banach dual of $E_a^0$ sample by sample.
If desired, its restriction to $C_c^\infty(\R^6)$ has a random
distribution version.  One direct construction expands locally in
Fourier functions multiplied by a compact cutoff.  Their variances
are uniformly bounded by \eqref{eq:target-field-bound} with $a=0$;
the expected squared local $H^{-s}$ norm is therefore finite for
$s>3$.  A countable exhaustion by cubes yields consistent local
distributions.  No path-space assertion is needed here.
\end{proof}
